\section{Heterogeneidad de las metas de gestión: ¿requiere el diseño vigente un tratamiento diferenciado?}

La sección anterior mostró que el marco regulatorio vigente reúne bajo la denominación de metas de gestión indicadores que miden objetos distintos. Algunas metas se refieren a resultados observados en la prestación, como la continuidad, la presión o la cobertura; otras verifican capacidades o instrumentos de gestión, como la actualización de catastros; y otras evalúan la ejecución física de inversiones, actividades o medidas de mejora, o la utilización de los recursos asignados.

La experiencia internacional examinada en el Anexo~\ref{anexo:experiencia_internacional} ofrece un punto de comparación para evaluar esta heterogeneidad. Los cinco casos revisados ---Italia, Inglaterra y Gales, Portugal, Chile y Colombia--- fueron seleccionados para combinar cercanía institucional y diversidad de modelos regulatorios, sin pretender una revisión exhaustiva. Aunque difieren en la forma de fijar metas y evaluar a los prestadores, comparten, con distintos alcances, una separación entre la evaluación de las condiciones observadas en la prestación y el control de las inversiones o actividades destinadas a producirlas. Asimismo, las distintas dimensiones evaluadas conservan una identificación propia, aun cuando algunas generen consecuencias económicas. Esta evidencia no determina el diseño aplicable al caso peruano; su utilidad es comparativa. La justificación de cualquier ajuste debe descansar en la correspondencia entre el objeto medido, los incentivos que genera el instrumento y la conducta que se pretende disciplinar.

La cuestión previa a la construcción del ICG es qué debe entenderse por desempeño regulatorio y qué indicadores resultan adecuados para medirlo. Si el índice busca sintetizar los resultados alcanzados por la empresa prestadora, la incorporación de variables referidas a los medios utilizados para producirlos puede mezclar dimensiones conceptualmente distintas y dificultar su interpretación. Sobre esta base, corresponde examinar la conveniencia de sustituir el actual esquema de metas de gestión por indicadores de desempeño referidos a resultados de cobertura y calidad del servicio, cuya agregación se efectuaría mediante el ICG, y reservar para las inversiones, actividades y demás medios de gestión instrumentos específicos de seguimiento y fiscalización.

La discusión parte de la delimitación de los indicadores que deberían representar el desempeño y de las dimensiones que, bajo este enfoque, correspondería resumir mediante el ICG. A continuación, examina la función del índice tras su desvinculación de los incrementos tarifarios y los problemas asociados con la utilización del ICI y del ICG como mecanismos de activación sancionadora, frente a una alternativa basada en la fiscalización de obligaciones específicas. Finalmente, analiza el eventual reconocimiento del sobrecumplimiento y ordena las alternativas mediante un principio de separación de instrumentos que sirve de enlace con la formalización de la Sección~4.

\subsection{Objeto del ICG: delimitación de las metas de desempeño}

Para delimitar el contenido del ICG es necesario precisar qué objeto regulatorio se busca resumir mediante el índice. Si este pretende representar el desempeño de la prestación, sus componentes deberían corresponder principalmente a condiciones observables del servicio. Las inversiones y los demás medios utilizados para producir esos resultados requieren, en cambio, mecanismos de seguimiento y verificación acordes con su naturaleza. La experiencia internacional ofrece un referente para esta separación: pese a las diferencias entre los marcos examinados, la evaluación de los resultados y el control de las inversiones se realizan mediante instrumentos diferenciados.\footnote{Las consecuencias asociadas al desempeño varían entre jurisdicciones e incluyen ajustes económicos y mecanismos reputacionales; cuando corresponde, el incumplimiento de obligaciones exigibles puede dar lugar además a sanciones.} La Tabla~\ref{tab:comparacion_internacional_separacion} resume esta evidencia y el Anexo~\ref{anexo:experiencia_internacional} desarrolla el detalle por jurisdicción.

\begin{table*}[!t]
\centering
\caption{Separación entre la evaluación del desempeño, sus consecuencias y el control de las inversiones en los marcos examinados}
\label{tab:comparacion_internacional_separacion}
\scriptsize
\renewcommand{\arraystretch}{1.16}
\setlength{\tabcolsep}{3.8pt}
\begin{tabularx}{\textwidth}{
>{\raggedright\arraybackslash}p{1.65cm}
>{\raggedright\arraybackslash}p{3.45cm}
>{\raggedright\arraybackslash}X
>{\raggedright\arraybackslash}p{3.55cm}}
\hline
\textbf{Jurisdicción} &
\textbf{Evaluación del desempeño} &
\textbf{Consecuencia del incumplimiento del desempeño} &
\textbf{Tratamiento de las inversiones} \\
\hline
Italia &
Macroindicadores con clases iniciales y objetivos de mantenimiento o mejora. &
Premios o penalidades por macroindicador, según la clase y la trayectoria aplicable. &
Programación tarifaria y penalidad separada por incumplimiento de la planificación de inversiones. \\
Inglaterra y Gales &
Compromisos comunes o específicos por empresa, con niveles y trayectorias verificables. &
Consecuencias financieras por compromiso mediante \emph{outcome delivery incentives}. &
Los \emph{price control deliverables} verifican productos, hitos, fechas y puesta en operación. \\
Portugal &
Evaluación multidimensional mediante referencias y calificaciones comparativas por entidad. &
La calificación es principalmente informativa y reputacional; las obligaciones vinculantes se controlan por separado. &
Planes, contratos y decisiones tarifarias controlan las obligaciones concretas fuera de la calificación ordinaria. \\
Chile &
Ranking normalizado de desempeño relativo. &
La posición en el ranking no genera por sí sola una sanción; las infracciones se determinan separadamente. &
Los planes de desarrollo y las obras comprometidas se fiscalizan fuera del ranking anual. \\
Colombia &
Estándares generales con metas anuales y gradualidades según el prestador. &
Determinadas metas generan descuentos en componentes tarifarios, sin sustituir otros mecanismos de control. &
El Plan de Obras e Inversiones Regulado controla activos planeados, ejecutados y puestos en operación. \\
\hline
\end{tabularx}
\begin{minipage}{\textwidth}
\footnotesize
\textit{Nota:} La tabla resume la arquitectura predominante de cada caso y no reproduce la totalidad de sus indicadores, obligaciones ni consecuencias económicas o sancionadoras. La versión ampliada se presenta en la Tabla~\ref{tab:anexo_experiencia_internacional} del Anexo~\ref{anexo:experiencia_internacional}. \textit{Fuente:} Elaboración propia sobre la base de la revisión internacional desarrollada en dicho anexo.
\end{minipage}
\end{table*}

\paragraph{Resultados y medios de implementación.}

Los indicadores utilizados para evaluar el desempeño se concentran principalmente en dimensiones de acceso y en condiciones observadas durante la prestación, como continuidad, presión, calidad, interrupciones, tratamiento o atención a los usuarios. En el régimen peruano, estas dimensiones se ubican principalmente en los aspectos de cobertura y calidad del servicio del SIIGEPSS, lo que permite examinar una eventual delimitación del ICG sin trasladar automáticamente al índice todas las dimensiones utilizadas para supervisar la gestión de las empresas prestadoras.

La correspondencia, sin embargo, no es automática. Como se observa en la Tabla~\ref{tab:grupos_indicadores_ep}, los aspectos de cobertura y calidad del servicio pueden incluir, junto con indicadores de resultado, variables que verifican la ejecución de intervenciones, como la instalación de conexiones o la renovación de redes. Por ello, la clasificación relevante debe atender al objeto directamente medido por cada indicador y no únicamente a su ubicación normativa. La Figura~\ref{fig:heterogeneidad_cobertura_calidad} ilustra esta diferencia.

\begin{figure}[!t]
\centering
\caption{Resultados e intervenciones dentro de los aspectos de cobertura y calidad del servicio}
\label{fig:heterogeneidad_cobertura_calidad}
\begin{tikzpicture}[
font=\scriptsize,
aspecto/.style={
draw,
rounded corners=2pt,
align=center,
text width=2.5cm,
minimum height=0.8cm,
inner sep=4pt
},
objeto/.style={
draw,
rounded corners=2pt,
align=center,
text width=3.7cm,
minimum height=0.8cm,
inner sep=4pt
},
conexion/.style={
-{Latex[length=1.7mm,width=1.2mm]},
line width=0.6pt
},
encabezado/.style={
font=\footnotesize\bfseries,
align=center
}
]
\node[encabezado] at (0,1.15) {Aspecto normativo};
\node[encabezado] at (4.6,1.15) {Objeto predominantemente medido};
\node[aspecto] (cobertura) at (0,0.15) {Cobertura};
\node[objeto] (resultado_cob) at (4.6,0.55) {Resultados de acceso y cobertura};
\node[objeto] (ejecucion_cob) at (4.6,-0.45) {Ejecución física de conexiones u otras intervenciones};
\draw[conexion] ([yshift=2mm]cobertura.east) to[out=0,in=180] (resultado_cob.west);
\draw[conexion] ([yshift=-2mm]cobertura.east) to[out=0,in=180] (ejecucion_cob.west);
\node[aspecto] (calidad) at (0,-2.45) {Calidad del servicio};
\node[objeto] (resultado_cal) at (4.6,-2.05) {Resultados de calidad y continuidad del servicio};
\node[objeto] (ejecucion_cal) at (4.6,-3.05) {Ejecución física o funcional de intervenciones};
\draw[conexion] ([yshift=2mm]calidad.east) to[out=0,in=180] (resultado_cal.west);
\draw[conexion] ([yshift=-2mm]calidad.east) to[out=0,in=180] (ejecucion_cal.west);
\end{tikzpicture}
\begin{minipage}{0.96\columnwidth}
\footnotesize
\textit{Nota:} La figura distingue el objeto predominantemente medido por los indicadores de cobertura y calidad del servicio; su ubicación normativa no determina por sí sola su incorporación al ICG. \textit{Fuente:} Elaboración propia sobre la base del Sistema de Indicadores e Índices de la Gestión de los Prestadores de los Servicios de Saneamiento.
\end{minipage}
\end{figure}

La distinción central es entre \emph{resultados} y \emph{medios de implementación}. Los primeros describen una condición alcanzada en el acceso o en la calidad de la prestación; los segundos verifican los instrumentos, productos o acciones utilizados para producirla. En adelante, este estudio denomina \emph{metas de desempeño} a las metas expresadas mediante indicadores que reflejan directamente esos resultados. Una eventual delimitación del ICG se evaluaría, por ello, a partir del objeto medido por cada indicador, mientras que las inversiones, actividades, capacidades de gestión o utilización de recursos conservarían los mecanismos de seguimiento y control que correspondan a su naturaleza.

Para mantener continuidad con la formalización posterior, se denomina \(J^{S}\) al conjunto de indicadores que miden directamente resultados observables de cobertura y calidad del servicio. Asimismo, se anticipa la distinción entre \(R^{S}\), que comprende las inversiones vinculadas con esos resultados; \(R^{G}\), que reúne las intervenciones asociadas con otras metas de gestión que cuentan con un indicador propio; y \(R^{F}\), que comprende las inversiones y medidas de mejora cuyo impacto no puede evaluarse mediante otro indicador aplicable. En \(A\) y \(A'\), estas últimas permanecen sujetas al avance financiero focalizado; desde \(A''\), se verifican mediante indicadores específicos de cumplimiento físico o funcional agrupados en \(J^{F}\). La definición formal de estos conjuntos y de su relación se presenta en la Subsección~4.1.

Esta clasificación no agota el universo de indicadores que quedan fuera de \(J^{S}\). Otras dimensiones regulatorias pueden conservar instrumentos propios sin corresponder a una intervención discreta comprendida en los grupos anteriores. Del mismo modo, la ejecución financiera no constituye un objeto regulatorio adicional, sino una forma de verificar determinadas intervenciones.

La distinción entre resultados y medios tiene una implicancia directa para la interpretación del cumplimiento. Una inversión, la adquisición de un equipo, la instalación de conexiones o la renovación de una red puede contribuir a mejorar el servicio, pero su ejecución no equivale al resultado que busca producir. La cadena de resultados utilizada en gestión y evaluación distingue precisamente entre insumos, actividades, productos, resultados e impactos \citep{oecd2023glossary}. Además, la ejecución de una intervención, su entrada en operación y la materialización de sus efectos pueden ocurrir en momentos distintos. Por ello, la ejecución anual no debe interpretarse mecánicamente como equivalente al desempeño anual observado.

La relación de estos objetos con las decisiones de la empresa también puede diferir. Los resultados responden a decisiones operativas y gerenciales, pero pueden verse afectados por condiciones estructurales de la prestación. La ejecución de inversiones, por su parte, puede depender de etapas del ciclo de inversión, procesos de contratación, disponibilidad de recursos, autorizaciones administrativas o actuaciones de terceros. En el sector peruano de saneamiento, \citet{escobar2025determinantes} identifican factores gerenciales, institucionales, financieros y geográficos asociados con la ejecución de inversiones de las empresas prestadoras. La experiencia chilena ofrece una manifestación institucional de esta separación: una posición desfavorable en el ranking de calidad no constituye por sí sola una infracción, cuya determinación requiere un procedimiento distinto (Anexo~\ref{anexo:chile}).

Los demás aspectos comprendidos en el SIIGEPSS seguirían cumpliendo funciones regulatorias relevantes aunque no necesariamente incidieran en el ICG. Los indicadores de sostenibilidad y solvencia económico-financiera pueden servir para seguimiento y alerta; los de eficiencia empresarial, para evaluar dimensiones operativas o comerciales; los de inversiones y utilización de reservas, para verificar recursos programados; y los de gobernanza, cumplimiento normativo o focalización del subsidio, para los fines específicos que justifican su existencia. Esta diferenciación de finalidades encuentra respaldo en \citet{behn2003measure}, quien advierte que las medidas de desempeño responden a usos distintos y que una misma medida no necesariamente es adecuada para todos ellos.

Dos indicadores que actualmente pueden incidir en el cálculo del ICG ilustran el problema con particular claridad: la relación de trabajo y el porcentaje de agua no facturada (ANF). Ninguno mide directamente una condición observada de la prestación ---como continuidad, presión, calidad o cobertura---, sino una variable intermedia asociada con la eficiencia económica u operativa de la empresa. Con el criterio de clasificación utilizado en este estudio, ambos corresponden a indicadores de eficiencia empresarial y no a las metas de desempeño comprendidas en \(J^{S}\), con independencia de su incidencia actual en el índice compuesto.

La incorporación de estas variables también afecta la interpretación del ICG. El mecanismo de dilución desarrollado formalmente en la Sección~4 para la ejecución financiera no depende de la naturaleza particular de dicha variable. En términos de agregación, cualquier componente adicional que compita por ponderación dentro del promedio reduce, manteniendo lo demás constante, la sensibilidad del índice frente a los resultados de cobertura y calidad que se pretende resumir. Por ejemplo, si un índice simple contiene inicialmente dos metas de resultado, cada una recibe un peso de \(1/2\); si se incorpora un indicador intermedio adicional, el peso individual de cada resultado se reduce a \(1/3\), aun cuando no haya cambiado la importancia sustantiva de las condiciones que estos representan.

El vínculo con los resultados comprendidos en \(J^{S}\) no es, sin embargo, igual en ambos casos. La relación de trabajo es un ratio que vincula costos operativos e ingresos y cuya mejora no representa por sí misma una condición directamente experimentada por el usuario; su clasificación como variable de eficiencia empresarial es, por ello, directa. El ANF presenta un vínculo productivo más próximo con determinadas condiciones del servicio, porque las pérdidas comerciales y físicas pueden afectar la disponibilidad efectiva del recurso y, a través de ella, resultados como continuidad o presión. Esa relación causal no lo convierte, sin embargo, en el resultado final observado: constituye una condición operativa intermedia que contribuye a producirlo. Su evolución, además, depende de un conjunto continuo de decisiones operativas, comerciales y de inversión, por lo que tampoco corresponde tratarlo como una intervención discreta dentro de \(R^{S}\), \(R^{G}\) o \(R^{F}\).

\paragraph{Delimitación de las metas de desempeño.}

La delimitación anterior identifica, dentro de la categoría más amplia de metas de gestión, las variables que podrían integrar un índice destinado a resumir resultados directamente observables de acceso o de prestación. La expresión \emph{metas de desempeño} resulta preferible a \emph{metas de calidad}, pues comprende también dimensiones de cobertura y otras condiciones cuantificables del servicio. El marco colombiano vigente desde 2026 ofrece un paralelo al combinar estándares de cobertura, continuidad, calidad y satisfacción del usuario mediante metas anuales con gradualidades diferenciadas (Anexo~\ref{anexo:colombia}).

La selección de los componentes del ICG debería partir del concepto regulatorio que se pretende representar. \citet{OECD2008CompositeIndicators} señalan que un indicador compuesto requiere un marco conceptual explícito que oriente la selección y combinación de variables. En el sector de agua potable, esta elección puede modificar materialmente la evaluación obtenida: \citet{PicazoTadeo2008Quality} muestran que incorporar explícitamente la calidad del servicio en la medición de la eficiencia de las empresas de agua produce resultados significativamente distintos de aquellos obtenidos mediante medidas basadas únicamente en cantidades. Esta exigencia se extiende a la ponderación y agregación: \citet{greco2019methodological} muestran que las reglas utilizadas pueden introducir compensaciones entre dimensiones. Por ello, combinar resultados con variables de implementación o eficiencia dentro de un mismo promedio puede dificultar la interpretación del índice y permitir que la mejora de un componente compense el deterioro de otro sin que exista una correspondencia sustantiva entre ambos.

El argumento central no depende de que otras jurisdicciones hayan adoptado exactamente la misma arquitectura. La comparación internacional muestra que la separación es operativamente viable, pero la razón principal para evaluarla proviene de la lógica del propio instrumento. Si una medida agrega conjuntamente resultados y medios, puede orientar esfuerzo hacia variables que mejoran la métrica sin producir una mejora equivalente en el objetivo regulatorio y reducir la sensibilidad del índice frente a las condiciones que pretende hacer visibles. La literatura sobre incentivos multitarea refuerza esta preocupación: \citet{holmstrom1991multitask} muestran que los incentivos sobre dimensiones fácilmente medibles pueden desplazar esfuerzo desde otras actividades, mientras que \citet{baker1992incentive} explica cómo una medida imperfectamente alineada puede inducir acciones que mejoran el indicador sin mejorar en igual proporción el objetivo subyacente.\footnote{\citet{smith1995unintended} identifica problemas relacionados en la medición del desempeño del sector público, como visión de túnel, suboptimización y fijación en la medida.} La Sección~4 formaliza esta lógica para la ejecución financiera; el mismo mecanismo de ponderación proporciona una referencia conceptual para evaluar la inclusión de indicadores intermedios, como la relación de trabajo o el ANF, aunque estos no se incorporen de manera simétrica como variables de decisión dentro del modelo.

\paragraph{Alcance de la clasificación y de la formalización.}

La clasificación propuesta es analítica y no describe íntegramente el diseño vigente. En particular, indicadores de eficiencia empresarial como la relación de trabajo y el ANF pueden incidir actualmente en el cálculo del ICG, por lo que su tratamiento mediante mecanismos separados de reporte, supervisión y, cuando corresponda, fiscalización constituye una modificación evaluada y no una descripción del \emph{statu quo}. La Sección~4 desarrolla formalmente la clasificación de las intervenciones y representa explícitamente la relación de trabajo en \(A_0\), mientras que el ANF y las demás dimensiones para las que no se especifica una tecnología propia permanecen fuera del conjunto de decisiones modeladas. Esta acotación permite concentrar el análisis formal en los incentivos asociados con la ejecución de inversiones sin suponer que los demás indicadores carezcan de relevancia regulatoria.

La delimitación define el contenido potencial del ICG. La cuestión siguiente es qué función puede cumplir el índice una vez separado de los instrumentos destinados a verificar inversiones y otras obligaciones de implementación, y qué consecuencia corresponde al incumplimiento de una meta de desempeño.

\subsection{Función del ICG y tratamiento del incumplimiento de las metas de desempeño}

Delimitar el contenido del ICG no implicaría reducir la evaluación regulatoria al valor de un único índice. El resultado global y sus componentes podrían publicarse conjuntamente y, por separado, mantenerse los indicadores de sostenibilidad financiera, eficiencia empresarial, cumplimiento de inversiones, capacidades de gestión y demás dimensiones relevantes. El objetivo es que cada señal conserve una interpretación identificable y que las consecuencias asociadas con su incumplimiento respondan al objeto regulatorio correspondiente.

Dentro de este diseño, el ICG cumpliría principalmente una función de síntesis, información y señalización reputacional del desempeño. Esta función era compatible con su anterior vinculación a los incrementos tarifarios. Sin embargo, el régimen vigente ya no condiciona los incrementos tarifarios base al cumplimiento de las metas de gestión. Por ello, no resulta directamente trasladable al diseño peruano un mecanismo en el que el incumplimiento de una meta de desempeño produzca, por sí mismo, una devolución de ingresos tarifarios previamente reconocidos. Esta restricción institucional no elimina la posibilidad de asociar consecuencias al desempeño, pero modifica el instrumento mediante el cual estas pueden materializarse.

En esta arquitectura, el incumplimiento de una meta comprendida en \(J^{S}\) se reflejaría directamente en su indicador individual de cumplimiento y, a través de este, en el ICG. Para mantener la notación de la Sección~4, sea \(\bar q_j\) la meta de desempeño aplicable al indicador \(j\). En aquellos indicadores para los que un mayor valor representa mejor desempeño,
\[
q_j<\bar q_j
\quad\Rightarrow\quad
ICI_j<100.
\]
El menor \(ICI_j\) se reflejaría, a su vez, en una señal de desempeño desfavorable. En la formalización posterior esta medida se expresa en términos normalizados mediante \(c_j=ICI_j/100\). El efecto inmediato de no alcanzar la meta es, por tanto, una reducción del indicador individual y de la señal agregada, sin que ese resultado permita identificar por sí solo la causa que lo produjo.

La consecuencia de medición no agota, sin embargo, el tratamiento jurídico de la meta. Cuando una meta de desempeño constituya además una obligación específica cuyo incumplimiento genere una consecuencia pecuniaria, esta debería determinarse respecto de dicha obligación y no mediante el nivel agregado del ICG. Esta es la configuración que la Sección~4 representa mediante las consecuencias específicas asociadas con \(J^{S}\). De manera adicional, una misma dimensión puede encontrarse sujeta a un estándar normativo autónomamente exigible, cuya fiscalización conserva su propio fundamento.

La efectividad de la función reputacional depende de la forma y frecuencia con que los resultados sean difundidos. En la práctica, el alcance reputacional del ICG ha sido limitado. La evaluación del cumplimiento de las metas de gestión se difundía principalmente con ocasión de las revisiones tarifarias, cuyos horizontes suelen abarcar varios años, y sus resultados se incorporaban dentro de los respectivos estudios tarifarios.\footnote{En los estudios tarifarios publicados hasta 2022, esta información se presentaba habitualmente en la subsección ``Cumplimiento de metas de gestión del quinquenio regulatorio'', incluida dentro de la sección ``Perfil de la empresa''. Los primeros estudios tarifarios publicados en 2023 mantuvieron esta estructura; posteriormente, durante ese mismo año, comenzaron a publicarse estudios que ya no incorporaban dicha subsección, patrón que se mantiene en los estudios posteriores.} Esta modalidad de difusión no configura un sistema periódico, directamente accesible y comparable entre empresas. Por ello, en su configuración actual, el ICG difícilmente puede generar una presión reputacional semejante a la que produciría la publicación frecuente, sistemática y fácilmente consultable de resultados comparables.

Una difusión más regular y visible fortalecería esa función. La publicación periódica del ICG y de sus componentes permitiría observar la evolución de cada empresa, facilitar comparaciones entre prestadores e informar a usuarios, juntas de accionistas, municipalidades propietarias y otros actores. Esta posibilidad se aproxima a los mecanismos de \emph{benchmarking} y \emph{sunshine regulation} descritos por \citet{marques2010towards}, en los que la publicidad comparativa del desempeño genera presión reputacional sin que el resultado deba determinar necesariamente las tarifas. La información comparativa también serviría como insumo para mecanismos más exigentes de \emph{yardstick competition}; en el sentido de \citet{shleifer1985theory}, estos utilizan el desempeño o los costos de empresas comparables como referencia para la regulación.

La meta de desempeño debe distinguirse, a su vez, de un estándar normativo autónomamente exigible. No alcanzar una trayectoria de mejora fijada como meta de gestión no equivale necesariamente a infringir una obligación normativa distinta de calidad. Una empresa puede obtener un resultado inferior a la meta y, por esa razón, registrar un menor \(ICI_j\), sin que de ello se desprenda automáticamente el incumplimiento de otro estándar aplicable a la misma dimensión. En cambio, si el resultado supone además la vulneración de una obligación autónoma de continuidad, presión, calidad u otra condición de prestación, la consecuencia correspondiente se determinaría mediante el régimen aplicable a esa obligación.

En términos conceptuales, se distinguen así dos referencias cuando ambas existen para una misma dimensión: \(\bar q_j\), que representa la meta de desempeño utilizada posteriormente en la Sección~4, y \(q_j^{N}\), que representa un estándar normativo autónomamente exigible. El incumplimiento de \(\bar q_j\) reduce el indicador de cumplimiento y puede, cuando la meta constituya una obligación específica sancionable, generar la consecuencia correspondiente sin depender del ICG agregado. El incumplimiento de \(q_j^{N}\) se evalúa separadamente conforme al régimen normativo aplicable. La Sección~4 formaliza la primera referencia y abstrae de la segunda como umbral adicional.

La diversidad de consecuencias observada internacionalmente muestra que no existe una única respuesta regulatoria frente al incumplimiento del desempeño. La opción evaluada para el Perú responde, por tanto, a la configuración del marco vigente y no a la reproducción de un mecanismo extranjero específico. El elemento comparado relevante es que las distintas obligaciones conservan una identificación propia y pueden sujetarse a instrumentos diferenciados.

Separar del ICG las inversiones y demás obligaciones de implementación exige determinar cómo deben verificarse y qué conducta debe activar una eventual consecuencia sancionadora. Esta cuestión no puede resolverse únicamente modificando el contenido del índice, pues el régimen vigente utiliza también los ICI y el ICG para definir determinados supuestos de infracción. Esa arquitectura sancionadora presenta problemas propios que conviene examinar primero.

\subsection{Problemas de la activación sancionadora mediante el ICI y el ICG}

Una eventual redefinición del ICG requeriría revisar también su correspondencia con el régimen sancionador. En el diseño vigente, la consecuencia sancionadora puede activarse a partir del ICG o de los ICI, aunque estos índices agregan o resumen obligaciones de distinta naturaleza. A su vez, la base económica de la multa puede construirse a partir de inversiones o medidas de mejora cuya ejecución habría sido postergada. Esta arquitectura plantea tres cuestiones relacionadas: la función de los umbrales frente a la determinación de responsabilidad, la posible incidencia múltiple de una misma intervención dentro del sistema de medición y la correspondencia entre la obligación cuyo cumplimiento se pretende asegurar, la variable utilizada para medirla y la conducta que puede generar una consecuencia sancionadora.

\textit{Primera cuestión: umbrales e imputabilidad.} La función económica del régimen sancionador es generar incentivos para el cumplimiento de las obligaciones regulatorias. El principio de razonabilidad del Reglamento busca que incumplir no resulte más ventajoso que cumplir y considera el beneficio ilícito y la probabilidad de detección entre los criterios relevantes. El beneficio ilícito puede comprender costos evitados o postergados derivados de gastos o inversiones que debieron ejecutarse oportunamente. La tesis económica que sustenta esta metodología es que la sanción esperada debe neutralizar la ventaja derivada del incumplimiento, idea desarrollada a partir de \citet{becker1968crime} y sistematizada posteriormente por \citet{stigler1970optimum}, \citet{polinsky2000economic} y \citet{polinsky2007theory}.

Los umbrales previstos para el ICI y el ICG modifican esta lógica al determinar qué incumplimientos pueden activar las tipificaciones correspondientes. Aunque la obligación sustantiva consista en alcanzar íntegramente una meta o ejecutar una intervención, el incumplimiento solo queda comprendido en estas tipificaciones cuando el ICI cae por debajo del \(80\%\) o el ICG por debajo del \(85\%\). Se generan así márgenes dentro de los cuales puede subsistir un incumplimiento parcial sin que se active la consecuencia sancionadora: una meta puede permanecer parcialmente incumplida mientras su ICI sea igual o superior al \(80\%\), y el incumplimiento agregado no activa la tipificación mientras el ICG alcance al menos el \(85\%\). En estos casos, la empresa puede conservar parte de la ventaja asociada con la porción de la obligación pendiente sin que opere el mecanismo sancionador destinado a neutralizarla.

El umbral aplicado al ICG genera un problema complementario. Como este índice se obtiene mediante el promedio de los ICI, un mayor nivel de cumplimiento en unas metas puede compensar el menor cumplimiento de otras y evitar que se active la tipificación agregada. Si, por ejemplo, dos metas registran ICI de \(90\%\) y \(82\%\), ninguna alcanza el cumplimiento íntegro, pero el ICG asciende a \(86\%\). El promedio supera entonces el umbral de \(85\%\) y ambos ICI permanecen también por encima del umbral individual de \(80\%\), por lo que los incumplimientos observados no activarían ninguna de estas tipificaciones. El esquema admite, por tanto, dos formas relacionadas de tolerancia: una respecto del cumplimiento individual de cada meta y otra derivada de la compensación entre metas dentro del índice agregado. Esta última permite que el nivel relativamente mayor de cumplimiento de una obligación contribuya a desactivar la consecuencia asociada con el menor cumplimiento de otra.

Estos efectos se producen antes del análisis de responsabilidad dentro del procedimiento administrativo sancionador. El numeral 30.2 del Reglamento General de Fiscalización y Sanción contempla causales individualizadas de exención, entre ellas el caso fortuito o fuerza mayor debidamente comprobados, el hecho determinante de tercero, la orden obligatoria de autoridad competente, el error inducido por la Administración y la subsanación voluntaria antes del inicio del procedimiento. Asimismo, establece que el cumplimiento parcial o defectuoso de las normas u obligaciones que configuran una infracción no exime por sí mismo de responsabilidad. Cuando el índice no cruza el umbral previsto en la tipificación, el incumplimiento queda fuera de ese supuesto sancionador antes de que corresponda examinar las circunstancias que podrían determinar o excluir la responsabilidad del administrado. El punto de corte opera, por ello, con independencia de las causas del incumplimiento y alcanza tanto a situaciones que podrían ser responsabilidad de la empresa como a aquellas que podrían encontrarse amparadas por una causal de exención.

Podría sostenerse que estos márgenes de tolerancia responden a las limitaciones de gestión que enfrentan algunas empresas prestadoras. Tales restricciones pueden ser relevantes para el diseño regulatorio, pero no justifican por sí solas un punto de corte sancionador uniforme. Si determinadas empresas requieren una trayectoria gradual de mejora, esa circunstancia puede reflejarse en la fijación de las metas y en sus cronogramas de cumplimiento. El umbral opera de manera distinta: una vez establecida la obligación como exigible, permite que subsista un incumplimiento parcial por el solo hecho de que el índice permanezca por encima del punto de corte y, en el caso del ICG, permite además que el cumplimiento relativamente mayor de otras obligaciones contribuya a evitar la activación de la tipificación agregada.

La discusión no implica que todo incumplimiento parcial deba generar automáticamente una sanción. La responsabilidad requiere una conducta previamente tipificada y que, atendiendo a las circunstancias del caso, pueda ser imputada al administrado. El problema identificado es más específico: los umbrales pueden debilitar el incentivo al cumplimiento íntegro al excluir determinados incumplimientos antes de que se evalúen esas circunstancias y, mediante el ICG, permitir que el menor cumplimiento de unas obligaciones se compense con un mayor nivel de cumplimiento de otras. Esto plantea la necesidad de examinar si el uso del ICI y del ICG como puntos de activación guarda correspondencia con la obligación cuyo cumplimiento se pretende asegurar.

\textit{Segunda cuestión: incidencia múltiple y doble contabilización en la medición.} Una misma intervención puede incidir en más de una medida de cumplimiento. Como se señaló en la Sección~2, el EP-57-A reduce una fuente de este problema hacia adelante, aunque no la elimina en estudios tarifarios aprobados bajo diseños anteriores. Mientras subsistan metas de avance financiero de alcance amplio, una misma inversión puede continuar formando parte del EP-57 y, simultáneamente, estar vinculada con otra meta específica.

En estos casos existe una duplicación directa dentro del sistema de medición. Una intervención puede incidir en el ICI del avance financiero por el gasto asociado con su ejecución y, al mismo tiempo, en el ICI de otra meta que verifica su ejecución física o el cumplimiento de una condición específica. Si ambos indicadores forman parte del ICG, la misma inversión adquiere más de una incidencia dentro de la medida agregada. El problema puede persistir, por tanto, mientras permanezcan vigentes estudios tarifarios aprobados bajo la configuración anterior.

La incidencia múltiple puede presentarse también cuando una misma inversión está vinculada con más de una meta de resultado. La instalación de \textit{data loggers} ilustra este supuesto: los equipos pueden formar parte de una intervención de inversión y, al mismo tiempo, resultar necesarios para obtener o acreditar los registros utilizados en la evaluación de continuidad y presión. Si la ausencia de esos equipos afecta simultáneamente los ICI de ambas metas, una única inversión incide más de una vez en la evaluación del cumplimiento y, cuando ambos indicadores integran el ICG, también en la medida agregada. A diferencia del caso anterior, la duplicación no proviene de combinar una meta de avance financiero con otra meta específica, sino de la vinculación de una misma inversión con varios resultados evaluados separadamente.

Una variante puede presentarse con la relación de trabajo. A diferencia de las metas asociadas con inversiones específicas, este indicador refleja una relación agregada entre costos operativos desembolsables e ingresos operacionales, cuya senda prevista puede depender de un conjunto amplio de inversiones y acciones de gestión. En los casos revisados, la base económica atribuida a su incumplimiento comprende una fracción importante del programa de inversiones y puede superponerse con inversiones vinculadas con otras metas específicas. En EMSAPUNO S.A., por ejemplo, la base de inversión atribuida a la relación de trabajo ascendió a aproximadamente S/ 10,76 millones, frente a S/ 13,57 millones del programa financiado con recursos internamente generados; en EPS EMSAPA CALCA S.A. alcanzó aproximadamente S/ 807 mil. Como contraste, en EPS Moyobamba la base atribuida al ANF se limitó a inversiones específicas por S/ 13,5 mil en equipos de presión y macromedidores y en la repotenciación y calibración del banco de medidores. Estos casos muestran que la magnitud de la superposición puede variar considerablemente según la meta y la forma en que se vinculan las inversiones.

El problema puede trasladarse posteriormente a la determinación de la sanción. Si el valor de una misma inversión no ejecutada se incorpora en la base económica asociada con el incumplimiento de continuidad y vuelve a incorporarse en la correspondiente a presión, se produce una doble valorización del costo evitado o postergado. Lo mismo puede ocurrir si una inversión incluida en la base amplia asociada con la relación de trabajo vuelve a ser considerada en la base de una meta específica. En ambos casos, una única inversión es contabilizada más de una vez al determinar la consecuencia económica. Este problema es distinto de la doble contabilización en el ICG, aunque ambos pueden originarse en la misma arquitectura de vinculación entre inversiones y metas.

La eventual imposición de sanciones autónomas plantea una cuestión adicional. Si las distintas metas afectadas por una misma omisión dan lugar a más de una sanción, corresponde examinar la aplicación del principio de \emph{non bis in idem}, cuya configuración exige identidad de sujeto, hecho y fundamento, así como las reglas aplicables al concurso de infracciones. La vinculación de una misma inversión con dos metas no permite afirmar por sí sola la existencia de \emph{non bis in idem}, pero sí exige verificar si las consecuencias recaen sobre conductas y fundamentos realmente distintos. Con independencia de ese análisis, la base económica no debería reiterar el valor de una misma inversión.

Conforme con la distinción planteada en la Subsección~3.1, cuando una intervención de \(R^{S}\) contribuye a uno o más resultados de \(J^{S}\), su vinculación con varias metas no justifica que reciba varias incidencias por el mismo hecho ni que su valor económico se reutilice en las bases asociadas con distintos incumplimientos. En la formalización de la Sección~4, la posible participación de una intervención en más de un componente físico activo de \(M^m\) se recoge mediante \(\delta_r^m\).

\textit{Tercera cuestión: correspondencia entre el objeto de cumplimiento y su medición.} El tercer problema aparece cuando la variable utilizada para evaluar una inversión no coincide con el objeto cuyo cumplimiento se pretende asegurar. Este es el caso de las metas de avance financiero. Si la obligación consiste en ejecutar un proyecto dentro del plazo y conforme a las condiciones previstas, medir su cumplimiento mediante el porcentaje del presupuesto gastado puede registrar como incumplimiento una diferencia que no proviene de una menor ejecución del proyecto, sino de haberlo completado a un costo inferior al programado.

El caso de EPS SEDA HUÁNUCO S.A. ilustra esta situación. El proyecto ``Descolmatación y Limpieza del río Higueras aguas arriba de la Captación Canchan'' fue completado por S/ 38\,000 frente a los S/ 101\,695 programados. Pese a que el alcance y la finalidad fueron reconocidos como cumplidos, el menor gasto redujo el ICI financiero a \(13.79\%\); si se hubiera reconocido el presupuesto íntegro de la intervención, el ICI habría ascendido a \(36.03\%\). El ahorro, sin embargo, fue excluido del beneficio ilícito utilizado para determinar la multa.\footnote{La Resolución de Dirección de Sanciones N.\textdegree~329-2025-SUNASS-DS reconoció el cumplimiento del alcance de la intervención y excluyó el ahorro del beneficio ilícito. La Resolución de Gerencia General N.\textdegree~290-2025-SUNASS-GG confirmó este criterio. El ICI permaneció por debajo del umbral de \(80\%\) porque otras dos intervenciones asociadas con la meta no fueron acreditadas, por lo que el reconocimiento del ahorro no modificó la multa final de 11.92 UIT.}

El caso muestra una asimetría entre la medición del cumplimiento y la determinación de la sanción. El ahorro puede reducir el ICI y, a través de este, deteriorar la señal de desempeño de la empresa o contribuir a activar una consecuencia sancionadora, aun cuando posteriormente ese mismo ahorro no sea considerado beneficio ilícito para calcular la multa. La exclusión del ahorro de la base económica evita que se le atribuya directamente una ventaja ilícita, pero no corrige la señal de incumplimiento generada previamente por el indicador financiero.

La diferencia resulta relevante incluso cuando la consecuencia final sea una amonestación. Si el proyecto de inversión fue ejecutado dentro del plazo y conforme a las condiciones previstas, y se encuentra operativo, un costo inferior al programado no representa por sí mismo una menor ejecución de la obligación. Puede reflejar una mayor eficiencia en la ejecución, siempre que el menor gasto no provenga de una reducción del alcance, la calidad, la capacidad, la vida útil u otra condición establecida para el proyecto. Medir el cumplimiento de la inversión mediante el gasto puede, por tanto, asociar una consecuencia reputacional o sancionadora con una diferencia financiera que no refleja necesariamente un incumplimiento de la intervención.

Los tres problemas responden a dimensiones distintas del diseño, pero convergen en una misma cuestión: la correspondencia entre la obligación que se pretende asegurar, la forma utilizada para medir su cumplimiento y la conducta que puede generar una consecuencia sancionadora. Esta correspondencia proporciona el punto de partida para evaluar una fiscalización y tipificación vinculadas directamente con las obligaciones específicas.

Los problemas antes señalados se sintetizan en la Tabla~\ref{tab:cuestiones_ici_icg}.

\begin{table}[t]
\centering
\caption{Síntesis de las cuestiones identificadas}
\label{tab:cuestiones_ici_icg}
\small
\renewcommand{\arraystretch}{1.15}
\begin{tabularx}{\columnwidth}{
>{\raggedright\arraybackslash}p{0.32\columnwidth}
>{\raggedright\arraybackslash}X}
\hline
\textbf{Cuestión} & \textbf{Síntesis} \\
\hline
Umbrales e imputabilidad &
Los umbrales del ICI y del ICG pueden excluir incumplimientos parciales antes de evaluar su imputabilidad y permitir compensaciones entre metas. \\
\hline
Incidencia múltiple y doble contabilización &
Una misma inversión puede incidir en varias metas o reutilizarse en las bases económicas asociadas con distintos incumplimientos. \\
\hline
Correspondencia entre objeto y medición &
La variable utilizada para medir el cumplimiento puede no coincidir con el objeto de la obligación, como ocurre al aproximar la ejecución mediante el gasto. \\
\hline
\end{tabularx}
\end{table}

\subsection{Fiscalización y tipificación por obligación específica}

Una alternativa transversal a los problemas anteriores consiste en desplazar la unidad de análisis sancionadora desde el nivel alcanzado por el ICI o el ICG hacia la obligación específica cuyo cumplimiento se pretende asegurar. El índice conservaría su función de evaluación del desempeño, mientras que la determinación de un incumplimiento sancionable requeriría identificar directamente la conducta u obligación exigible. De este modo, el cumplimiento relativamente mayor de otras metas no podría neutralizar el incumplimiento de una obligación distinta mediante el promedio agregado.

\paragraph{Verificación directa de las obligaciones de implementación.}

En el caso de los proyectos de inversión, la fiscalización verificaría directamente que la intervención o sus componentes hayan sido ejecutados dentro de los plazos y conforme a las condiciones previstas, así como su entrada efectiva en operación cuando esta forme parte de la obligación. El gasto realizado continuaría siendo información relevante para el seguimiento financiero y para determinar, cuando corresponda, la ventaja económica asociada con la parte no ejecutada, pero no sustituiría la comprobación del proyecto. Esta separación permite distinguir la ejecución física, la entrada en operación y el costo incurrido, dimensiones que pueden evolucionar de manera diferente.

Los \emph{price control deliverables} de Ofwat y el control de los planes de desarrollo en Chile ofrecen antecedentes de esta forma de verificación. En ambos casos, el seguimiento de las intervenciones puede atender a productos, cantidades, hitos, plazos o puesta en operación, sin identificar el cumplimiento de la inversión únicamente con el porcentaje de recursos gastados (Anexo~\ref{anexo:experiencia_internacional}).

Las demás obligaciones de implementación seguirían una lógica equivalente. Las intervenciones de \(R^{G}\) pueden verificarse mediante los productos, hitos, condiciones o plazos definidos para la meta correspondiente. Para las intervenciones de \(R^{F}\), la secuencia evaluada mantiene inicialmente la verificación financiera focalizada y, desde \(A''\), la sustituye por indicadores específicos de cumplimiento físico o funcional agrupados en \(J^{F}\). El instrumento de fiscalización debería corresponder, en cada caso, con aquello que efectivamente se exige, en lugar de trasladar su cumplimiento a un índice destinado a resumir resultados del servicio.

Para las inversiones de \(R^{S}\), la distinción anterior implica que el resultado observado no permite presumir por sí solo el cumplimiento o incumplimiento de la intervención asociada. Su ejecución debe verificarse directamente conforme al alcance, plazo y demás condiciones exigibles. La Tabla~\ref{tab:resultado_inversion_consecuencias} muestra las cuatro combinaciones posibles entre el cumplimiento del resultado y el de la inversión.

\begin{table*}[!t]
\centering
\caption{Tratamiento conjunto del resultado de desempeño y de la inversión asociada}
\label{tab:resultado_inversion_consecuencias}
\small
\renewcommand{\arraystretch}{1.20}
\setlength{\tabcolsep}{5pt}
\begin{tabularx}{\textwidth}{
>{\raggedright\arraybackslash}p{2.7cm}
>{\raggedright\arraybackslash}p{2.7cm}
>{\raggedright\arraybackslash}X}
\hline
\textbf{Resultado \(J^{S}\)} &
\textbf{Inversión \(R^{S}\)} &
\textbf{Tratamiento regulatorio en la alternativa evaluada} \\
\hline
Cumple &
Cumple &
El indicador de desempeño refleja el cumplimiento de la meta y no existe incumplimiento de la obligación de inversión. \\
No cumple &
Cumple &
El incumplimiento se refleja en un menor \(ICI_j\) y en la señal reputacional. Cuando la meta de desempeño constituya además una obligación específica sancionable, puede generar la consecuencia correspondiente con independencia del ICG agregado. No existe incumplimiento de la inversión, pues esta fue ejecutada conforme a la obligación exigible. \\
Cumple &
No cumple &
El resultado favorable se refleja en el \(ICI_j\), pero no elimina el incumplimiento de la obligación de inversión cuando esta sea autónomamente exigible. La consecuencia correspondiente se determina respecto de esa obligación. \\
No cumple &
No cumple &
El menor \(ICI_j\) refleja el incumplimiento de la meta de desempeño y puede generar, cuando corresponda, su consecuencia específica. Separadamente, la inversión no ejecutada puede generar la consecuencia asociada con esa obligación. Ambas respuestas recaen sobre objetos regulatorios distintos. \\
\hline
\end{tabularx}
\begin{minipage}{\textwidth}
\footnotesize
\textit{Nota:} La tabla supone que la inversión constituye una obligación autónomamente exigible y que su eventual consecuencia sancionadora depende del régimen aplicable. El incumplimiento de la meta de desempeño no se utiliza para presumir el incumplimiento de la inversión. \textit{Fuente:} Elaboración propia.
\end{minipage}
\end{table*}

\paragraph{Correspondencia sancionadora y base económica.}

A partir de esta separación, la consecuencia sancionadora y su base económica deberían vincularse con la obligación efectivamente incumplida. Cuando una misma inversión contribuye a varias metas de desempeño, reconstruir la base económica de la sanción a partir de cada una de ellas puede llevar a incorporar reiteradamente el mismo costo evitado o postergado. Si la infracción se vincula directamente con la inversión incumplida, el valor económico de la parte no ejecutada se asociaría con esa obligación y se contabilizaría una sola vez, con independencia del número de resultados a los que contribuya.

Esta duplicación económica debe distinguirse del principio de \emph{non bis in idem}. Una misma inversión puede contribuir a diferentes resultados y estos, a su vez, pueden corresponder a obligaciones autónomas de calidad del servicio. La existencia de varias consecuencias no implica por sí sola una vulneración de dicho principio, cuya aplicación requiere examinar la identidad de sujeto, hecho y fundamento y las reglas sobre concurso de infracciones. La separación propuesta exige precisamente identificar la conducta y el fundamento de cada consecuencia, evitando además que el valor de una misma inversión se reutilice en varias bases económicas.

La misma lógica permite separar el cumplimiento del proyecto del gasto realizado. Una vez ejecutada la intervención dentro del plazo y conforme a las condiciones previstas, un costo inferior al programado no reduciría por sí mismo su nivel de cumplimiento. El ahorro podría conservar relevancia para el seguimiento financiero o para las reglas aplicables al destino de los recursos remanentes, pero no debería registrarse como una porción físicamente no ejecutada cuando el alcance exigido haya sido completado. El método tarifario italiano ofrece un antecedente de separación entre las consecuencias asociadas con la planificación de inversiones y los premios o penalidades correspondientes al desempeño en calidad técnica (Anexo~\ref{anexo:italia}). El tratamiento del ahorro se formaliza en la Sección~4 mediante el monto reconocido una vez cumplidas las condiciones previstas para la intervención.

La Tabla~\ref{tab:comparacion_diseno_separacion} resume las principales diferencias entre el diseño vigente y la alternativa evaluada en materia de inversiones.

\begin{table*}[t]
\caption{Verificación y tratamiento sancionador de las inversiones en el diseño vigente y en la alternativa evaluada}
\label{tab:comparacion_diseno_separacion}
\centering
\small
\renewcommand{\arraystretch}{1.25}
\setlength{\tabcolsep}{5pt}
\begin{tabularx}{\textwidth}{
@{}
>{\raggedright\arraybackslash}p{0.19\textwidth}
>{\raggedright\arraybackslash}X
>{\raggedright\arraybackslash}X
@{}
}
\hline
\textbf{Elemento} &
\textbf{Diseño vigente} &
\textbf{Diseño alternativo evaluado} \\
\hline
Verificación de las inversiones &
La ejecución financiera puede utilizarse como medida de cumplimiento, aunque el gasto realizado no necesariamente coincida con la ejecución del proyecto en las condiciones previstas. &
Se verificaría directamente la ejecución de la intervención o de sus componentes dentro de los plazos y conforme a las condiciones previstas, así como su entrada en operación, registrando separadamente el costo realizado. \\
Incidencia en el ICG &
La ejecución de una inversión puede modificar directamente el ICG, incluso antes de que se materialice el resultado que busca producir. &
La inversión se verificaría separadamente y solo incidiría en el ICG de desempeño a través del resultado comprendido en \(J^{S}\) que contribuya a producir. \\
Unidad de análisis sancionadora &
La activación de determinadas tipificaciones depende del nivel alcanzado por el ICI o el ICG. &
La determinación del incumplimiento se realizaría respecto de la obligación específica, sin depender del cumplimiento agregado de otras metas. \\
Base económica &
Una misma inversión puede incidir en más de una meta y, si no se identifica separadamente, su valor puede incorporarse reiteradamente al determinar las consecuencias asociadas con distintos incumplimientos. &
El valor de la parte no ejecutada se asociaría con la obligación específica y se contabilizaría una sola vez. \\
Ahorro en la ejecución &
Un menor gasto puede reducir el cumplimiento de una meta financiera aun cuando el alcance de la intervención haya sido completado. &
Un costo inferior al programado no reduciría por sí mismo el cumplimiento de la inversión cuando se hayan satisfecho el alcance y las demás condiciones exigibles. \\
\hline
\end{tabularx}
\begin{flushleft}
\footnotesize
\textit{Nota:} La columna correspondiente al diseño alternativo resume los elementos vinculados con la verificación y el tratamiento sancionador de las obligaciones de implementación examinados en esta subsección. Sus propiedades económicas se formalizan en la Sección~4. \textit{Fuente:} Elaboración propia.
\end{flushleft}
\end{table*}

El tratamiento separado no elimina la función preventiva del seguimiento de las inversiones. Los cronogramas, procesos de contratación, hitos de ejecución y fechas de entrada en operación permiten identificar retrasos antes de que se configure el incumplimiento definitivo o antes de que sus efectos se reflejen en las condiciones del servicio. La cuestión no es, por tanto, si estas variables deben observarse, sino cuál es la función regulatoria que corresponde a cada una y qué hecho debe constituir la conducta sancionable.

El criterio de correspondencia puede extenderse también a capacidades, actividades e indicadores de eficiencia empresarial que no se representan mediante las intervenciones consideradas en el modelo. Su exclusión del ICG no supone que carezcan de relevancia ni que no puedan estar sujetos a obligaciones exigibles. Cuando exista una obligación autónoma y expresamente tipificada, su incumplimiento podrá generar la consecuencia correspondiente, pero la medición y la base económica de la sanción deberán guardar relación con la conducta que se pretende disciplinar.

Con ello se distinguen la evaluación del desempeño, la fiscalización de las obligaciones de implementación y la exigibilidad de estándares normativos autónomos. Queda pendiente, en cambio, un problema de signo contrario: el reconocimiento de los resultados que exceden las metas.

\subsection{Reconocimiento del sobrecumplimiento de las metas de desempeño}

Una vez delimitado el ICG a metas de desempeño y separadas las consecuencias asociadas con las obligaciones específicas, surge una cuestión adicional: cómo tratar los resultados que superan el nivel regulatorio establecido. El diseño vigente trunca cada ICI en \(100\%\). Esta regla evita que el sobrecumplimiento de una meta compense directamente el incumplimiento de otra dentro del cumplimiento ordinario, pero también elimina el incentivo marginal proveniente de ese componente del ICG una vez alcanzada la meta.

La experiencia internacional muestra que el desempeño superior puede recibir un reconocimiento específico, aunque los mecanismos utilizados difieren entre jurisdicciones. Italia incorpora premios dentro del sistema de incentivos asociado con la calidad técnica; Inglaterra y Gales contempla, para determinados compromisos, pagos por desempeño superior al nivel comprometido mediante los \emph{outcome delivery incentives}; y Portugal utiliza premios y sellos de calidad como mecanismos de reconocimiento reputacional separados de la evaluación ordinaria. Estos antecedentes no reproducen una arquitectura única ni determinan el mecanismo aplicable al caso peruano, pero muestran que el desempeño superior puede reconocerse mediante un instrumento diferenciado del utilizado para identificar el cumplimiento ordinario o el incumplimiento (Anexo~\ref{anexo:experiencia_internacional}).

La alternativa evaluada mantiene separado el cumplimiento ordinario del sobrecumplimiento. Para una meta \(\bar q_j\) respecto de la cual un mayor valor representa mejor desempeño, el componente ordinario conservaría la forma
\begin{equation}
ICI_j
=
100\min\left\{
\frac{q_j}{\bar q_j},1
\right\}.
\end{equation}
De este modo, cuando \(q_j>\bar q_j\), el \(ICI_j\) ordinario no superaría el \(100\%\). El desempeño adicional se identificaría mediante un componente específico y acotado, sin modificar el cumplimiento base de las demás metas ni la configuración de sus consecuencias.

Esta distinción permite que, en la extensión \(B^{+}\) desarrollada en la Sección~4.4, el sobrecumplimiento reconocido se incorpore de manera explícita a un ICG extendido sin confundirse con \(C_S\), que mantiene el cumplimiento ordinario truncado en uno. Como desde \(B\) las consecuencias específicas ya no dependen del nivel agregado del ICG, un resultado superior en una dimensión no desactiva el incumplimiento ni la consecuencia correspondiente a otra. El componente adicional cumple así una función reputacional y de reconocimiento sin reintroducir el umbral agregado como mecanismo de compensación sancionadora.

El reconocimiento adicional solo tendría sentido cuando la mejora sea verificable, sostenible, socialmente valiosa y razonablemente atribuible a las decisiones de la empresa. No todas las metas satisfacen estas condiciones. Algunas presentan límites técnicos, problemas de medición, dificultades para distinguir el efecto de las decisiones empresariales de factores externos, riesgos de manipulación o beneficios marginales decrecientes a partir de determinado nivel. En consecuencia, cualquier mecanismo de reconocimiento requeriría definir ex ante los indicadores elegibles, el nivel a partir del cual comienza a reconocerse el desempeño adicional, un límite máximo y las circunstancias que impedirían su otorgamiento.

Por ejemplo, si la meta de continuidad fuese de \(20\) horas por día y la empresa alcanzara \(22\), el ICI ordinario permanecería en \(100\%\). Las dos horas adicionales ---equivalentes a un sobrecumplimiento relativo de \(10\%\) respecto de la meta--- podrían incorporarse al componente adicional de \(B^{+}\) dentro del rango reconocido. Si otra meta registrara un ICI inferior a \(100\%\), dicho sobrecumplimiento no eliminaría ese incumplimiento ni impediría la aplicación de la consecuencia específica que correspondiera.

El reconocimiento puede adoptar una forma reputacional, regulatoria o económica. La publicación de resultados destacados, rankings específicos, premios o sellos constituye la alternativa más directamente compatible con la función informativa del ICG. También podrían evaluarse otros mecanismos regulatorios cuando exista una relación suficientemente clara entre el desempeño superior y el beneficio generado. Los premios y sellos utilizados en Portugal ofrecen un antecedente de reconocimiento reputacional separado de la calificación ordinaria; los mecanismos de Italia y de Ofwat muestran, por su parte, que determinados diseños regulatorios pueden incorporar además incentivos económicos positivos asociados con resultados superiores a la referencia establecida (Anexo~\ref{anexo:experiencia_internacional}).

Una eventual bonificación económica en el Perú requeriría, sin embargo, un instrumento, una fuente de financiamiento y una habilitación propios. No podría interpretarse como una restitución de la antigua condicionalidad general de los incrementos tarifarios base, dado que estos ya no dependen del cumplimiento de las metas de gestión. Asimismo, el reconocimiento debería permanecer identificado separadamente del cumplimiento ordinario y no utilizarse para neutralizar consecuencias correspondientes a otras metas, obligaciones de implementación o estándares normativos autónomamente exigibles.

La Sección~4.4 formaliza esta posibilidad mediante \(B^{+}\), que mantiene el cumplimiento base \(C_S\), reconoce de manera acotada el sobrecumplimiento dentro de un ICG extendido y añade una bonificación predeterminada por los resultados adicionales verificables.

\subsection{Síntesis: principio de separación de instrumentos}

La discusión anterior permite formular un posible principio de diseño: cada objeto regulatorio debería vincularse con el instrumento que mejor corresponda a su función. La Tabla~\ref{tab:separacion_instrumentos} sintetiza esta correspondencia sin reproducir el tratamiento particular de las inversiones desarrollado en la Subsección~3.4.

\begin{table}[!t]
\centering
\caption{Correspondencia evaluada entre objetos regulatorios e instrumentos}
\label{tab:separacion_instrumentos}
\small
\renewcommand{\arraystretch}{1.20}
\setlength{\tabcolsep}{5pt}
\begin{tabularx}{\columnwidth}{@{}p{0.37\columnwidth}X@{}}
\hline
\textbf{Objeto regulatorio} & \textbf{Tratamiento regulatorio evaluado} \\
\hline
Metas de desempeño de \(J^{S}\) &
ICG de desempeño, transparencia y comparación entre empresas. Cuando exista una obligación específica sancionable, su consecuencia se determina sin depender de un umbral agregado del ICG. \\
Estándares normativos autónomamente exigibles &
Fiscalización y aplicación de las medidas previstas por la regulación correspondiente, con independencia del ICG. \\
Obligaciones de implementación &
Seguimiento y verificación según el objeto exigido ---ejecución, producto, hito, condición o puesta en operación--- y, cuando corresponda, aplicación de consecuencias respecto de la obligación específica. \\
Indicadores de sostenibilidad, solvencia, eficiencia y otros propósitos específicos &
Mecanismos especializados de reporte, seguimiento, supervisión o fiscalización según su finalidad, sin incidencia automática sobre el ICG de desempeño. \\
Sobrecumplimiento verificable &
Componente adicional y acotado en \(B^{+}\), separado del cumplimiento ordinario y sin capacidad para neutralizar incumplimientos de otras obligaciones. \\
\hline
\end{tabularx}
\begin{flushleft}
\footnotesize
\textit{Nota:} Las obligaciones de implementación comprenden las intervenciones de \(R^{S}\), \(R^{G}\) y \(R^{F}\), cuyo tratamiento específico se desarrolla en la Subsección~3.4 y se formaliza en la Sección~4. Los indicadores de eficiencia empresarial incluyen, entre otros, la relación de trabajo y el ANF. \textit{Fuente:} Elaboración propia.
\end{flushleft}
\end{table}

El principio preservaría la exigibilidad de las obligaciones sin exigir que todas incidan en una misma medida. El ICG serviría para sintetizar y comparar los resultados de la prestación; las obligaciones de implementación y los estándares normativos conservarían mecanismos propios de fiscalización; y las dimensiones de eficiencia, sostenibilidad u otros propósitos regulatorios se seguirían mediante instrumentos acordes con su función. El sobrecumplimiento, cuando resulte elegible, podría reconocerse de manera separada.

La arquitectura descrita puede alcanzarse mediante modificaciones incrementales del diseño de referencia. La Tabla~\ref{tab:secuencia_esquemas} presenta esta secuencia. \(A_0\) representa el diseño de referencia aplicado en la mayoría de los estudios tarifarios examinados; las transiciones posteriores aíslan progresivamente la focalización del avance financiero, la regla de activación de las consecuencias, la forma de verificación de \(R^{F}\), la delimitación del ICG y, finalmente, el reconocimiento del sobrecumplimiento.

\begin{table}[!t]
\centering
\caption{Secuencia de modificaciones evaluadas en la Sección~4}
\label{tab:secuencia_esquemas}
\small
\renewcommand{\arraystretch}{1.15}
\setlength{\tabcolsep}{4pt}
\begin{tabularx}{\columnwidth}{@{}p{0.18\columnwidth}X@{}}
\hline
\textbf{Transición} & \textbf{Modificación evaluada} \\
\hline
\(A_0\to A\) & Focaliza el avance financiero en \(R^{F}\) y retira la relación de trabajo del ICG y de las consecuencias pecuniarias. \\
\(A\to A'\) & Elimina el umbral agregado del ICG como condición de activación de las consecuencias pecuniarias. \\
\(A'\to A''\) & Sustituye la verificación financiera de \(R^{F}\) por indicadores específicos de cumplimiento físico o funcional de \(J^{F}\). \\
\(A''\to B\) & Delimita el ICG a los resultados comprendidos en \(J^{S}\). \\
\(B\to B^{+}\) & Añade un reconocimiento acotado del sobrecumplimiento y una bonificación por resultados adicionales verificables. \\
\hline
\end{tabularx}
\begin{flushleft}
\footnotesize
\textit{Fuente:} Elaboración propia.
\end{flushleft}
\end{table}

La secuencia permite evaluar cada modificación de manera incremental y evita presentar \(B^{+}\) como una propuesta indivisible. La Sección~4 formaliza el mecanismo de dilución asociado con la incorporación de variables distintas de los resultados, así como los cambios en los incentivos derivados de la focalización del avance financiero, la eliminación del umbral agregado y la sustitución de la verificación financiera. Su extensión a indicadores de eficiencia empresarial, como la relación de trabajo o el ANF, responde a la misma lógica de ponderación dentro de un índice compuesto, aunque estos no se incorporen de manera simétrica como variables explícitas de decisión en el modelo.